\documentclass[../../../include/open-logic-section]{subfiles}

\begin{document}

\olfileid{sth}{ord-arithmetic}{mult}

\olsection{序数乘法}

现在我们转向序数乘法，其处理方式与序数加法非常相似。假设我们要计算 $\alpha$ 乘以 $\beta$。为此，可以想象一个宽为 $\alpha$、高为 $\beta$ 的矩形网格；$\alpha$ 与 $\beta$ 的乘积，就是沿着每一行移动，然后移到下一行……直至遍历整个网格的结果。换言之，将 $\beta$ 中的\emph{每个}元素替换为 $\alpha$ 的一个\emph{副本}，就得到了 $\alpha$ 与 $\beta$ 的乘积。

为了将其形式化，我们直接对 $\alpha$ 和 $\beta$ 的笛卡尔积取逆字典序：

\begin{defn}
对于任意序数 $\alpha, \beta$，它们的乘积 $\alpha \ordtimes \beta = \ordtype{\alpha \times \beta, \rlexless}$。
\end{defn}
\noindent
我们必须再次确认这是一个良定义：

\begin{lem}\ollabel{ordtimeslessiswo}
对任意序数 $\alpha$ 和 $\beta$，$\tuple{\alpha \times \beta, \rlexless}$ 是一个良序。
\end{lem}

\begin{proof}
证明与 \olref[add]{ordsumlessiswo} 完全相同。
\end{proof}
\noindent
并且不难证明乘法满足以下性质：

\begin{lem}\ollabel{ordtimesrecursion}
对任意序数 $\alpha, \beta$：
\begin{align*}
	\alpha \ordtimes 0 &= 0\\
	\alpha \ordtimes (\beta \ordplus 1) &= 
		(\alpha \ordtimes \beta) \ordplus \alpha\\
	\alpha  \ordtimes \beta &= 
		\supstrict_{\delta < \beta}(\alpha \ordtimes \delta) && 
		\text{当 $\beta$ 是极限序数时}.
\end{align*}
\end{lem}

\begin{proof}
留作练习。
\end{proof}

事实上，正如加法的情形一样，我们本可以通过这些递归方程来定义序数乘法，而不是直接给出定义。同样与加法类似，乘法的某些性质也是我们所熟悉的：

\begin{lem}\ollabel{ordinalmultiplicationisnice}
若 $\alpha, \beta, \gamma$ 是序数，则：
\begin{enumerate}
	\item\ollabel{ordtimes1} 若 $\alpha \neq 0$ 且 $\beta < \gamma$，
	则 $\alpha \ordtimes \beta < \alpha \ordtimes \gamma$；
	\item\ollabel{ordtimes2} 若 $\alpha \neq 0$ 且 $\alpha \ordtimes
	\beta = \alpha\ordtimes\gamma$，则 $\beta = \gamma$；
	\item\ollabel{ordtimes3}  $\alpha \ordtimes (\beta \ordtimes
	\gamma) = (\alpha \ordtimes \beta) \ordtimes \gamma$；
	\item\ollabel{ordtimes4}  若 $\alpha \leq \beta$，则 $\alpha
	\ordtimes \gamma \leq \beta \ordtimes \gamma$；
	\item\ollabel{ordtimes5}  $\alpha \ordtimes (\beta \ordplus
	\gamma) = (\alpha \ordtimes \beta )\ordplus (\alpha\ordtimes
	\gamma)$。
\end{enumerate}
\end{lem}

\begin{proof}
留作练习。
\end{proof}

读者可按自己的兴趣去证明（或查阅）其他结果。但是，考虑到
\olref[ord-arithmetic][add]{ordsumnotcommute}，
以下结论应当不足为奇：

\begin{prop}
序数乘法不满足交换律：$2 \ordtimes \omega  =
\omega < \omega \ordtimes 2$。
\end{prop}

\begin{proof}
$2 \ordtimes \omega = \supstrict_{n < \omega}(2\ordtimes  n) = \omega \in \supstrict_{n < \omega}(\omega \ordplus n) = \omega \ordplus \omega = \omega \ordtimes 2$。
\end{proof}
\noindent 
同样，其直观解释也非常直接。为了计算 $2
\ordtimes \omega$，我们用两个对象替换每个自然数。如果简单地将所有“半数”插入自然数序列中，即考虑 $\Setabs{\nicefrac{n}{2}}{n \in \omega}$ 上的自然序，就会得到相同的序型。这样看来，该序型显然与 $\omega$ 本身的序型相同。但是，为了计算 $\omega \ordtimes 2$，则需要依次排下 $\omega$ 的两个副本，一个接在另一个之后。

\begin{prob}
证明
\olref[sth][ord-arithmetic][mult]{ordtimeslessiswo}、
\olref[sth][ord-arithmetic][mult]{ordtimesrecursion}、
以及
\olref[sth][ord-arithmetic][mult]{ordinalmultiplicationisnice}。
\end{prob}

\end{document}